\section{LoRA: adaptación de rango bajo}

\subsection{El dilema del ajuste fino con todos los parámetros}

Cuando deseamos adaptar un modelo de difusión preentrenado a un estilo específico o a un sujeto particular (como generar fotografías de una persona concreta), el método más directo es ajustar finamente todos los parámetros del modelo. Sin embargo, esto plantea problemas graves:

\begin{itemize}
    \item El U-Net de Stable Diffusion tiene aproximadamente 860M de parámetros; el ajuste fino con todos los parámetros requiere una gran cantidad de memoria gráfica (VRAM).
    \item Ajustar finamente un modelo grande con pocos datos (por ejemplo, 5--20 imágenes) provoca fácilmente \textbf{sobreajuste}.
    \item Cada versión personalizada necesita almacenar los pesos completos del modelo (varios GB), lo que incrementa el costo de despliegue.
    \item El ajuste fino puede dañar las capacidades generales aprendidas por el modelo preentrenado (\textbf{olvido catastrófico}).
\end{itemize}

\subsection{La idea central de LoRA}

LoRA (Low-Rank Adaptation) se basa en una hipótesis clave: \textbf{las actualizaciones de peso provocadas por el ajuste fino residen en un subespacio de rango bajo}.

Para una matriz de pesos $W_0 \in \mathbb{R}^{d \times k}$ dentro del modelo preentrenado, LoRA descompone la actualización en el producto de dos matrices de rango bajo:

$$
W = W_0 + \Delta W = W_0 + B \cdot A
$$

donde $B \in \mathbb{R}^{d \times r}$, $A \in \mathbb{R}^{r \times k}$ y $r \ll \min(d, k)$ es el rango.

\begin{importantbox}{Eficiencia de parámetros de LoRA}
Tomando como ejemplo una matriz de pesos de $1024 \times 1024$:
\begin{itemize}
    \item Ajuste fino con todos los parámetros: $1024 \times 1024 = 1,048,576$ parámetros entrenables.
    \item LoRA ($r=4$): $1024 \times 4 + 4 \times 1024 = 8,192$ parámetros entrenables.
    \item Reducción de parámetros: \textbf{128 veces}.
\end{itemize}
En los modelos de difusión, generalmente se aplica LoRA solo a las matrices $W_Q$, $W_K$, $W_V$, $W_O$ de las capas de atención, reduciendo aún más el número total de parámetros entrenables.
\end{importantbox}

\subsection{Aplicación de LoRA en modelos de difusión}

Estrategia de inicialización y entrenamiento de LoRA:

\begin{itemize}
    \item $A$ se inicializa con ruido gaussiano aleatorio, mientras que $B$ se inicializa como una matriz nula.
    \item Al comenzar el entrenamiento, $\Delta W = BA = 0$, garantizando que la salida inicial sea consistente con el modelo preentrenado.
    \item Solo se entrena $A$ y $B$, manteniendo congelados todos los parámetros originales $W_0$.
    \item Durante la inferencia, $\Delta W$ puede combinarse de nuevo en $W_0$, sin \textbf{aumentar ningún retraso de inferencia}.
\end{itemize}

Escenarios de uso típicos:
\begin{itemize}
    \item \textbf{Adaptación de estilo}: entrenar un LoRA con obras de un artista específico para generar imágenes en ese estilo.
    \item \textbf{Personalización de sujetos} (DreamBooth + LoRA): entrenar con varias fotos de un objeto o persona específica para generar nuevas imágenes de ese sujeto.
    \item \textbf{Inyección de conceptos}: enseñar al modelo a comprender nuevos conceptos visuales (como el logo de una marca específica).
\end{itemize}

\begin{warningbox}{Selección del rango de LoRA}
El rango $r$ es el hiperparámetro más crítico de LoRA:
\begin{itemize}
    \item Si $r$ es demasiado pequeño (por ejemplo, 1--2): la capacidad expresiva es insuficiente y el resultado del ajuste fino es pobre.
    \item Si $r$ es demasiado grande (por ejemplo, 128+): se acerca al ajuste fino con todos los parámetros, perdiendo la ventaja de eficiencia de parámetros.
    \item Recomendación práctica: para modelos de difusión, $r = 4 \sim 32$ suele ser una buena elección.
    \item Para tareas de adaptación de estilo, $r$ puede ser más pequeño (4--8); para tareas con condiciones complejas, $r$ necesita ser mayor (16--32).
\end{itemize}
\end{warningbox}

\begin{knowledgebox}{Combinabilidad de LoRA}
Una característica elegante de LoRA es su \textbf{combinabilidad}: múltiples módulos LoRA pueden combinarse dinámicamente durante la inferencia. Por ejemplo, un LoRA controla el estilo y otro controla el sujeto; ambos pueden cargarse simultáneamente y mezclar sus pesos para lograr "un artista específico pintando a una persona específica". Esto convierte a LoRA en la base de un ecosistema personalizado impulsado por la comunidad —existen miles de modelos LoRA en plataformas como Civitai—.
\end{knowledgebox}

\subsection{Resumen de la sección}

LoRA reduce drásticamente el número de parámetros y la carga computacional requerida para el ajuste fino mediante descomposición de rango bajo, haciendo posible adaptar modelos de difusión a gran escala con pocos datos en GPUs de nivel consumidor. Su estrategia de inicialización nula (similar en espíritu a ControlNet) garantiza que el ajuste fino no degrade las capacidades preentrenadas, mientras que la combinación de parámetros durante la inferencia elimina sobrecostes adicionales de latencia.

\newpage
